% Generated draft: requires mathematical and source review.
\documentclass[11pt]{article}
\usepackage[margin=25mm]{geometry}
\usepackage{fontspec}
\setmainfont{Latin Modern Roman}
\usepackage{amsmath,amssymb,mathtools,unicode-math}
\setmathfont{Latin Modern Math}
\usepackage{enumitem,needspace,xurl,hyperref}
\setlength{\parindent}{0pt}
\setlength{\parskip}{5pt}
\setlength{\emergencystretch}{3em}
\newcommand{\N}{\mathbb{N}}
\newcommand{\Z}{\mathbb{Z}}
\newcommand{\Q}{\mathbb{Q}}
\newcommand{\R}{\mathbb{R}}
\newcommand{\C}{\mathbb{C}}
\newcommand{\F}{\mathbb{F}}
\newcommand{\A}{\mathbb{A}}
\newcommand{\PP}{\mathbb{P}}
\newcommand{\E}{\mathbb{E}}
\newcommand{\eps}{\varepsilon}
\newcommand{\dd}{\,\mathrm{d}}
\begin{document}
\section*{UnsolvedMath ID: 20000729}
\textbf{Source problem code:} catalogue entry
\subsection{Definitions and mathematical statement}
Let \(\tau\in\mathbb H\), \(q=e^{2\pi i\tau}\), and \((a;q)_n=\prod_{j=0}^{n-1}(1-aq^j)\), with the convergent infinite-product convention for \(|q|<1\). Define from their unilateral sums
\[
 G(q)=\sum_{n\geq0}\frac{q^{n^2}}{(q;q)_n},\qquad
 H(q)=\sum_{n\geq0}\frac{q^{n^2+n}}{(q;q)_n}.
\]
The original motivation is the Rogers--Ramanujan identities
\[
 G(q)=\frac1{(q;q^5)_\infty(q^4;q^5)_\infty},\qquad
 H(q)=\frac1{(q^2;q^5)_\infty(q^3;q^5)_\infty}.
\]
It asks for a method proving modularity directly from the sum side, rather than taking these product evaluations as input. Normalize \(g(\tau)=q^{-1/60}G(q)\), \(h(\tau)=q^{11/60}H(q)\), with fractional powers interpreted through \(\tau\).

The aid's reported product-free route uses Bressoud's terminating Gaussian-polynomial identities, recurrence certification, and compact-open convergence to unary theta differences, followed by Poisson summation. The exact terminating identities and recurrences are not supplied in the record. A precise theta target for the standard normalization is
\[
 \Theta_r(\tau)=\sum_{n\equiv r\ (\mathrm{mod}\ 20)}q^{n^2/40},\qquad
 g=\frac{\Theta_1-\Theta_9}{\eta},\qquad
 h=\frac{\Theta_3-\Theta_7}{\eta},
\]
where \(\eta(\tau)=q^{1/24}\prod_{j\geq1}(1-q^j)\). The supplied route does not invoke either Rogers--Ramanujan product identity or Jacobi's triple product as an input. In this normalization its reported exact transformation laws are
\[
 \begin{pmatrix}g(-1/\tau)\\h(-1/\tau)\end{pmatrix}
 =\frac2{\sqrt5}\begin{pmatrix}
 \sin(2\pi/5)&\sin(\pi/5)\\
 \sin(\pi/5)&-\sin(2\pi/5)
 \end{pmatrix}\begin{pmatrix}g(\tau)\\h(\tau)\end{pmatrix},
\]
\[
 \begin{pmatrix}g(\tau+1)\\h(\tau+1)\end{pmatrix}
 =\begin{pmatrix}e^{-2\pi i/60}&0\\0&e^{22\pi i/60}\end{pmatrix}
 \begin{pmatrix}g(\tau)\\h(\tau)\end{pmatrix}.
\]
This specifies weight-zero vector-valued modularity, allowing the cusp pole indicated by \(q^{-1/60}\), rather than claiming both entries are holomorphic modular forms at every cusp.

The more general sufficient certificate in the aid has the following explicit finite form. For an integer \(M\geq1\), put
\[
 \Theta_{M,r}(\tau)=\sum_{n\equiv r\ (\mathrm{mod}\ 2M)}q^{n^2/(4M)}.
\]
After division by \(\eta\), the residue vector has finite transformation matrices
\[
 \mathsf S_{r,s}=(2M)^{-1/2}e^{-2\pi i rs/(2M)},\qquad
 \mathsf T_{r,s}=1_{r=s}e^{2\pi i(r^2/(4M)-1/24)}.
\]
If a finite coefficient space of linear combinations is preserved under the induced row actions of \(\mathsf S,\mathsf T\), and a recurrence-certified sum-side finitization converges compactly to those eta-denominated combinations, then it yields the corresponding vector-valued modular transformation laws. Relations such as \(\Theta_{M,r}=\Theta_{M,-r}\) must be respected, either by symmetric coefficient vectors or by taking the appropriate quotient. This is a sufficient certificate for sums admitting the required finite theta realization, not a universal algorithm deciding modularity of arbitrary q-hypergeometric series.
\subsection{Short English statement}
The source asks how to see modularity directly from q-hypergeometric sums. The supplied route passes through terminating Gaussian-polynomial identities and compact convergence to unary theta differences, then uses Poisson summation. For the Rogers--Ramanujan pair it gives explicit transformation matrices; finite theta-space closure is a sufficient certificate, not a general modularity decision procedure.
\subsection{Sources}
\begin{itemize}
\item Local registry: permanent UnsolvedMath ID \(20000729\), exact public code \texttt{catalogue entry}; \texttt{aim-analytic-number-theory-notes.json}, \texttt{notes[92]}.
\item AIM workshop \emph{Mock modular forms in combinatorics and arithmetic geometry}, March 8--12, 2010, Problem 1.2: \url{https://aimath.org/WWN/mockmodular/mockmodular.pdf}. This URL is supplied provenance; no PDF access is claimed.
\item Original motivation: the Rogers--Ramanujan identities for \(G,H\), without full bibliographic references.
\item Supplied research summary, ``A finite-theta certificate for sum-side modularity,'' attempt 1, cites Bressoud's terminating Gaussian-polynomial identities, Poisson summation, and finite Weil-matrix closure without full theorem citations. Original remarks and literature fields are empty.
\end{itemize}
\end{document}
